\documentclass[11pt]{letter}
\usepackage[utf8]{inputenc}
\usepackage{geometry}
\geometry{margin=1in}
\usepackage{hyperref}

\signature{Haythem Ghazouani, Ph.D.\\
Université de Tunis El Manar / Université de Carthage\\
\texttt{haythem.ghazouani@enicar.u-carthage.tn}}

\address{Haythem Ghazouani\\
Institut Supérieur d'Informatique\\
Université de Tunis El Manar\\
2 Rue Abou Rayhane Bayrouni\\
Ariana, 2080, Tunisia}

\date{\today}

\begin{document}

\begin{letter}{Editor-in-Chief\\
Swarm and Evolutionary Computation\\
Elsevier}

\opening{Dear Editor,}

We are pleased to submit our manuscript entitled \textbf{``Dual-Population Adversarial Strategy Adaptation (DPASA): A Meta-Optimization Framework for Global Optimization''} for consideration for publication in \textit{Swarm and Evolutionary Computation}.

\textbf{Summary of the Work:}
This paper introduces DPASA (Dual-Population Adversarial Strategy Adaptation), a novel population-based optimization algorithm that fundamentally decouples search learning from solution finding through a dual-population architecture. Unlike traditional metaheuristics that conflate search strategies with candidate solutions, DPASA maintains two co-evolving populations: one of \textit{strategy vectors} navigating the ``landscape of search behaviors'' and another of \textit{candidate solutions} navigating the fitness landscape. The algorithm incorporates five key innovations: (1) sinusoidal guidance functions for parsimonious trajectory encoding using only $4D$ parameters, (2) a negation operator---a deterministic anti-correlation mechanism for adversarial exploration in strategy parameter space, (3) success-history based parameter adaptation, (4) an elite archive for solution preservation, and (5) Nelder-Mead local search for precise convergence.

\textbf{Key Contributions:}
\begin{itemize}
    \item A novel dual-population meta-optimization architecture that enables simultaneous optimization of search strategy and solution vectors
    \item The negation operator as a structured adversarial mechanism that differs fundamentally from opposition-based learning by operating in strategy parameter space rather than solution space
    \item Comprehensive empirical validation on 15 classical benchmark functions, 29 CEC 2017 competition functions, and the CEC 2022 suite across multiple dimensions ($D \in \{10, 30, 50\}$)
    \item Comparison with CEC 2024 competition winners: DPASA achieves Friedman rank 2.00 on CEC 2017, ranking second to L-SHADE (1.97), with superior performance on specific deceptive multimodal landscapes
    \item Validation on six constrained engineering design problems with gaps of 0.11\%, 0.004\%, 0.01\%, 0.00\%, and 0.01\% on five problems, while honestly reporting structural limitations on the Pressure Vessel problem
\end{itemize}

\textbf{Why Swarm and Evolutionary Computation:}
We believe this manuscript is well-suited for \textit{Swarm and Evolutionary Computation} because it:
\begin{itemize}
    \item Introduces a genuinely novel algorithmic paradigm with clearly defined computational mechanisms (rather than superficial metaphorical rebranding)
    \item Provides rigorous empirical and statistical evaluation following best practices, including comparison with CEC 2024 competition winners
    \item Addresses the important challenge of exploration-exploitation balance through principled adversarial mechanisms
    \item Includes comparison with recent state-of-the-art methods (L-SRTDE, RDE, jSO, SHADE, L-SHADE, CMA-ES)
    \item Explicitly discusses algorithmic limitations and operational boundaries, providing guidance on when the algorithm should not be used
\end{itemize}

\textbf{Ethical Declarations:}
\begin{itemize}
    \item This manuscript has not been published elsewhere and is not under consideration by any other journal.
    \item All authors have approved the manuscript and agree with its submission.
    \item The research did not involve human participants or animals.
    \item The authors declare no conflicts of interest.
    \item No external funding was received for this research.
    \item Code and data are publicly available at: \url{https://github.com/haythemghz/DPASA-algorithm}
\end{itemize}

\textbf{Suggested Reviewers:}
We suggest the following experts who are familiar with the topics covered in this manuscript:
\begin{enumerate}
    \item Prof. Ponnuthurai N. Suganthan, Nanyang Technological University, Singapore (expertise: differential evolution, benchmark functions)
    \item Prof. Janez Brest, University of Maribor, Slovenia (expertise: adaptive DE variants, jSO algorithm)
    \item Prof. Andries Engelbrecht, Stellenbosch University, South Africa (expertise: swarm intelligence, fitness landscape analysis)
\end{enumerate}

We appreciate your consideration of this manuscript and look forward to your response.

\closing{Sincerely,}

\end{letter}
\end{document}
